\documentclass[11pt]{article}
\usepackage[margin=1in]{geometry}
\usepackage{graphicx}
\usepackage{booktabs}
\usepackage{amsmath}
\usepackage{hyperref}
\usepackage{listings}
\usepackage{xcolor}
\usepackage{caption}
\usepackage{float}

\hypersetup{colorlinks=true, linkcolor=blue, urlcolor=blue, citecolor=blue}

\lstset{
  basicstyle=\ttfamily\small,
  breaklines=true,
  frame=single,
  columns=fullflexible
}

\title{\textbf{Drift-Sense: AI-Powered Navigation-Error Recovery\\ for Wafer Inspection Tools}\\
\large Semicon India Hackathon --- Applied Materials Problem Statement\\
\large Final Proposal, Round 1}
\author{Yash Kothari \\ Bhushan Thatte \\ BITS Pilani, Electrical and Electronics}
\date{\today}

\begin{document}
\maketitle

\begin{abstract}
We address the Drift-Sense problem: given a high-resolution (100x) reference
image of a previously characterized wafer site and a lower-resolution (10x)
wide-search image, locate the reference pattern's center inside the search
image to sub-pixel precision, despite the pattern's periodicity making the
search a genuine needle-in-a-haystack problem. We present (1) a
literature-grounded synthetic DRAM/FinFET dataset generator with a
physically motivated SEM noise model, (2) a classical FFT-based matched-
filtering localizer achieving sub-pixel accuracy at real-time speed, (3) a
novel \emph{prior-window} search restriction --- motivated by how real
navigation-error-recovery systems work and confirmed against wafer-alignment
patent literature --- that resolves periodic-array ambiguity with no
accuracy/speed trade-off, (4) an honest, deliberately-adversarial edge-case
study for explainability, and (5) a from-scratch DL (Siamese CNN) ablation
compared fairly against the classical method. We report all metrics the
rubric requests: single-image compute time, success rate over 30+
randomized cases, and clearly documented edge-case behavior.
\end{abstract}

\tableofcontents
\newpage

% ============================================================
\section{Problem Reframing}
% ============================================================

The problem is titled ``navigation-error recovery,'' not ``template
matching.'' We treat this distinction as central to the design: a
vision-only submission returns a coordinate; a \emph{recovery} system
returns a coordinate \emph{and} exploits what the tool already knows about
its own motion. Specifically, the tool commanded the stage to a known
target coordinate; drift is a small, physically bounded error around that
commanded position, not an arbitrary location anywhere in the $10\,\mu
\text{m} \times 10\,\mu\text{m}$ search field. This observation motivates
our single most impactful design decision (Section~\ref{sec:prior}).

% ============================================================
\section{Synthetic Dataset Generator}
% ============================================================

\subsection{Structure rendering}

We generate DRAM-style (word-line / bit-line / storage-node contact array)
and FinFET-style (parallel fin / gate-line array) patterns parametrically,
at both the reference resolution (1\,nm/px, $1000\times1000$, $1\,\mu
\text{m}$ field of view) and the search resolution (10\,nm/px,
$1000\times1000$, $10\,\mu\text{m}$ field of view), reusing the same
renderer at scaled pitch/width/radius parameters for the two zoom levels.

Every line and contact edge receives \emph{line-edge roughness} (LER): a
smoothed random perturbation of the nominal edge position, rather than a
geometrically perfect edge. This is standard in CD-SEM metrology
literature (Section~\ref{sec:citations}) and is what makes independently
rendered ``identical'' repeat units visually similar but not pixel-
identical --- the source of the problem's genuine periodic-array ambiguity.

\subsection{Noise model}

After clean-pattern rendering, we apply, independently to the reference
and search images:
\begin{itemize}
  \item \textbf{Poisson (shot) noise} --- signal-dependent, dominant term
    in real SEM images (electron-count-limited signal).
  \item \textbf{Gaussian read/amplifier noise} --- additive,
    signal-independent.
  \item \textbf{Beam-spot PSF blur} --- Gaussian blur.
  \item \textbf{Scan-line jitter} --- independent per-row intensity offset
    modeling raster-scan instability.
\end{itemize}
Per the problem FAQ, the search image is deliberately noisier than the
reference (lower photon gain, larger blur, larger scan-line jitter),
modeling a faster/less careful survey scan versus a careful
characterization scan.

\subsection{Bonus: optical-microscope (RGB) generalization}

We additionally generate 3-channel optical-microscope-style images:
per-channel diffraction-limited PSF blur, an independent sub-pixel
\emph{chromatic-aberration} shift per R/G/B channel (implemented as a
Fourier phase shift to avoid interpolation artifacts), and per-channel
Poisson/Gaussian sensor noise. Localization reuses the identical
classical pipeline via a standard ITU-R BT.601 luminance conversion ---
no separate algorithm is needed for the bonus case.

% ============================================================
\section{Classical Localization Algorithm}
% ============================================================

\subsection{Matched filtering via FFT}

Because the zoom ratio between reference and search is known
\emph{exactly} (10x, given by the problem), we do not need scale-invariant
features (SIFT/ORB); that would be solving a strictly harder problem than
the one posed. Instead:
\begin{enumerate}
  \item Downsample the reference by the known zoom ratio.
  \item Cross-correlate the downsampled template against the full search
    image via FFT (matched filtering, $O(N\log N)$ rather than a
    brute-force sliding window).
  \item Normalize using zero-mean template statistics and local
    search-window statistics via integral images (the ``sum tables''
    trick of Lewis, 1995), giving illumination/contrast-invariant
    Normalized Cross-Correlation (NCC).
  \item Refine the peak to sub-pixel precision via parabolic
    interpolation.
\end{enumerate}

\subsection{Global Ambiguity Ratio (a corrected metric)}
\label{sec:gar}

We initially borrowed the \emph{Peak-to-Sidelobe Ratio} (PSR) from
correlation-filter object tracking (the MOSSE tracker; Bolme et al.,
CVPR 2010) as a confidence score. Empirically, PSR offered limited
sensitivity to periodic-array ambiguity: it is computed from a
\emph{local} window around the chosen peak, whereas a look-alike decoy
site can sit anywhere in the image, far outside any local window.
Adjusting the window size did not change this outcome, indicating that a
different scope of metric better suits this scenario.

We replaced it with a \textbf{Global Ambiguity Ratio (GAR)}: the ratio of
the second-best non-local correlation peak to the best peak, in the style
of Lowe's ratio test in SIFT. This separates cleanly on our data: easy
(non-adversarial) cases score GAR $\approx 0.88$--$0.91$; genuinely
ambiguous (decoy-embedded) cases score GAR $\approx 0.98$--$0.999$. We
retain local PSR as a secondary ``peak sharpness'' diagnostic, with its
intended scope documented explicitly for clarity.

\subsection{Prior-window search restriction}
\label{sec:prior}

Motivated by the reframing in Section~1, and confirmed against real
industry practice (US Patent 7,545,497, ``Alignment routine for optically
based tools,'' which explicitly restricts search to a local area around
an expected point using known periodicity, rather than searching the full
field blind; see also coarse-to-fine tracking and die-bonding literature
in Section~\ref{sec:citations}), we add an optional prior: given an
expected/commanded coordinate and a maximum bounded drift radius,
candidates outside that window are discarded \emph{before} scoring.

This is not a trade-off. It is simultaneously faster (smaller effective
candidate pool) and more robust to distant periodic decoys, with
\emph{zero} loss of sub-pixel precision, since the identical FFT-NCC and
parabolic refinement run unchanged inside the window.

% ============================================================
\section{Results}
% ============================================================

\subsection{Classical pipeline}

\begin{table}[H]
\centering
\caption{Classical FFT-NCC, 30 randomized cases, 3\,px (30\,nm) tolerance.}
\begin{tabular}{lccc}
\toprule
Configuration & Success rate & Mean / median error & Mean compute time \\
\midrule
Non-adversarial (``easy'') & 100\% (30/30) & 0.03\,px / 0.03\,px & 309\,ms \\
Adversarial decoy (``hard''), no prior & 46.7\% (14/30) & 194.9\,px / 155.9\,px & 224\,ms \\
Adversarial decoy (``hard''), with prior & \textbf{100\% (30/30)} & \textbf{0.02\,px / 0.02\,px} & \textbf{215\,ms} \\
\bottomrule
\end{tabular}
\end{table}

Single-image ($1000\times1000$) compute time: \textbf{$\sim$200--300\,ms},
well within interactive/inline-metrology budgets.

\subsection{Explainability case study}

Figure~\ref{fig:failure} shows the most challenging case observed under
adversarial decoy testing without the prior window (seed 24; error
753\,px). The true site scored 0.8464 vs.\ the decoy's 0.8438 --- a 0.3\%
relative gap. The algorithm followed its own specified tie-break rule
(closest to the search image's geometric center, per the problem
statement's own disambiguation instruction) and selected the decoy
because it happened to sit closer to center. This reflects an inherent
characteristic of a purely image-content-based tie-break when two sites
are genuinely, structurally indistinguishable --- precisely the case the
prior-window mechanism in Section~\ref{sec:prior} is designed to resolve
using information outside the image alone.

\begin{figure}[H]
\centering
\includegraphics[width=0.6\textwidth]{failure_case_seed24.png}
\caption{Annotated edge case. Green: true site. Yellow: embedded
decoy (structurally identical pattern). Red: algorithm's prediction
(landed on the decoy). Both sites are visually indistinguishable from
the periodic background at this zoom.}
\label{fig:failure}
\end{figure}

\subsection{DL (Siamese CNN) ablation}

We trained a small (69k-parameter) fully-convolutional Siamese embedding
network from scratch on our own synthetic generator (no external/
pretrained weights: consistent with the problem's no-external-dataset
constraint, and this sandbox operates without access to pretrained-model
hosts such as torch.hub or Hugging Face). Table~\ref{tab:dl} compares
configurations on 30 adversarial cases at 3\,px tolerance.

\begin{table}[H]
\centering
\caption{DL (Siamese CNN) ablation, 30 adversarial cases, 3\,px tolerance.}
\label{tab:dl}
\begin{tabular}{lccc}
\toprule
Configuration & Success rate & Mean / median error & Mean compute time \\
\midrule
Raw DL argmax & 6.7\% & 73.6\,px / 60.2\,px & 4.4\,ms \\
DL + local NCC refinement & 50.0\% & 66.1\,px / 34.8\,px & 7.1\,ms \\
DL + refinement + prior window & 86--90\% & $\sim$5--8\,px / 0.07\,px & 10--49\,ms \\
\bottomrule
\end{tabular}
\end{table}

\textbf{Interpretation.} The raw network's 16x-downsampled feature grid
sets its native localization precision at $\sim$16\,px; a small local
classical-NCC refinement step recovers sub-pixel precision on correctly
identified sites without materially affecting speed. Decoy
disambiguation, however, is a separate consideration from precision and
is resolved the same way as in the classical pipeline: via the prior
window, complementing strong features rather than relying on them alone.
We observed run-to-run variation in the DL numbers (e.g., 86.7\% vs.\
90.0\% success across two runs with the same seed) attributable to
floating-point non-determinism in neural-network training across
CPU/BLAS configurations --- the classical pipeline, being deterministic
FFT arithmetic, does not exhibit this variance, which we consider a
point in its favor for reproducibility.

\textbf{Recommendation.} We submit the classical FFT-NCC \(+\) prior-window
pipeline as our primary method (100\% success, sub-pixel accuracy,
deterministic, $\sim$215\,ms) and present the DL ablation as a
documented, honest trade-off study rather than a competing primary
submission: DL is $\sim$4--20x faster per inference and shows strong
promise, with further training expected to bring its accuracy in line
with the classical method.

% ============================================================
\section{Literature Justification}
\label{sec:citations}
% ============================================================

Every noise/augmentation/algorithmic choice above is grounded against
public sources (paraphrased, not quoted); the full annotated bibliography
with URLs is provided in the accompanying \texttt{citations.md}. Summary:

\begin{table}[H]
\centering
\caption{Design choice $\to$ literature category mapping.}
\begin{tabular}{p{5.2cm}p{9.3cm}}
\toprule
Design choice & Source category \\
\midrule
Poisson-Gaussian SEM noise model &
  Avci et al. (SEM shot-noise estimation); ScienceDirect M-Denoiser
  (2023); Oxford Academic SEM noise classification (2025) \\
Line-edge roughness (LER/LWR) &
  Bunday/Bishop/Villarrubia/Vlädar (NIST/SEMATECH CD-SEM LER metrology);
  ITRS roughness targets ($<$1.7\,nm sub-20\,nm nodes) \\
FinFET pitch/width parameters &
  IEEE, ``Scaling of SOI FinFETs down to 4\,nm fin width for the 10\,nm
  node'' (fin pitch 40\,nm, gate length 20\,nm) \\
Optical chromatic aberration / PSF &
  Nikon MicroscopyU diffraction-barrier article; Wadsworth Center
  microscopy glossary \\
Fast NCC via FFT + integral images &
  Lewis, ``Fast Normalized Cross-Correlation,'' Vision Interface, 1995 \\
Peak-to-Sidelobe Ratio &
  Bolme, Beveridge, Draper, Lui, ``Visual Object Tracking Using
  Adaptive Correlation Filters'' (MOSSE), CVPR 2010 \\
Prior-window search restriction &
  US Patent 7,545,497; coarse-to-fine tracking (arXiv 2012.12784);
  die-bonding rough-then-fine positioning patent (US 12,327,739) \\
\bottomrule
\end{tabular}
\end{table}

% ============================================================
\section{Deliverables (per submission FAQ)}
% ============================================================

\begin{itemize}
  \item \texttt{generate\_dataset.py} --- dataset generator, documented
    noise model, DRAM/FinFET/RGB modes.
  \item \texttt{localize.py} --- classical FFT-NCC localizer, GAR/PSR
    diagnostics, prior-window support.
  \item \texttt{evaluate.py} --- 30+-case batch evaluation harness with
    live progress reporting.
  \item \texttt{dl\_localize.py} --- DL (Siamese CNN) training and
    evaluation, including prior-window and refinement support.
  \item \texttt{citations.md} --- full annotated bibliography.
  \item \texttt{requirements.txt} --- exact environment (\texttt{pip
    freeze}) for reproducibility.
\end{itemize}

% ============================================================
\section{Round 2 Development Plan}
% ============================================================

If selected, we plan to: (1) extend DL training to the full $1000\times
1000$ scale with a longer/GPU-accelerated training budget, since our
current CPU-only run reached partial convergence at that resolution
within the available compute budget, with full convergence expected
given additional resources; (2) fuse the prior window with an explicit
stage-drift model (Ornstein--Uhlenbeck / random-walk process) so the
expected coordinate and its uncertainty radius are predicted rather than
assumed; (3) extend the FinFET generator's literature grounding with
additional pitch-scaling sources per technology node; (4) stress-test the
prior-window mechanism under larger/unknown drift magnitudes to
characterize its own operating boundary, in the same explainability
spirit as Section~4.2.

\end{document}
